% Source: https://artofproblemsolving.com/community/c3h1481123_graphing_question
% Stable identity: aops-topic-1481123-q01
For $x>0$ the power $x^{b}$ is positive. If $b>0$ it increases from $0$ (as $x\to0^+$) to $\infty$ (as $x\to\infty$), passing through $(1,1)$. If $b<0$, then $x^{b}=\dfrac{1}{x^{|b|}}$ decreases from $\infty$ to $0$, again through $(1,1)$. Multiplying by $a>0$ stretches the graph vertically (it now passes through $(1,a)$); multiplying by $a<0$ also reflects it in the $x$-axis.
\begin{enumerate}
\item[(a)] $a>0$, $b<0$: in the first quadrant, \textbf{decreasing}; $y\to\infty$ as $x\to0^+$ (the $y$-axis is an asymptote) and $y\to0$ as $x\to\infty$ (the $x$-axis is an asymptote). Like $y=\dfrac{1}{x}$.
\item[(b)] $a<0$, $b<0$: the reflection of (a) in the $x$-axis, in the fourth quadrant, \textbf{increasing}; $y\to-\infty$ as $x\to0^+$ and $y\to0$ as $x\to\infty$.
\item[(c)] $a>0$, $b>0$: in the first quadrant, \textbf{increasing}; $y\to0$ as $x\to0^+$ and $y\to\infty$ as $x\to\infty$. The curve bends upward (like $y=x^2$) if $b>1$, is a straight line through the origin if $b=1$, and bends downward (like $y=\sqrt x$) if $0<b<1$.
\item[(d)] $a<0$, $b>0$: the reflection of (c), in the fourth quadrant, \textbf{decreasing}; $y\to0$ as $x\to0^+$ and $y\to-\infty$ as $x\to\infty$.
\end{enumerate}
A table of values is a good way to start, but the two facts above (the sign of $a$ and whether $x^{b}$ grows or shrinks) already fix the shape. Typical examples with $a=\pm1$ (in (c) and (d) the solid curves have $b=2$ and the dashed ones $b=\dfrac{1}{2}$):

\begin{asy}
usepackage("amsmath");
size(8cm);
defaultpen(fontsize(10pt));
import graph;
real f1(real x){return 1/x;} real f2(real x){return -1/x;}
real g1(real x){return x^2;} real g2(real x){return sqrt(x);}
real g3(real x){return -x^2;} real g4(real x){return -sqrt(x);}
void myax(pair o){draw(o+(-0.2,0)--o+(3.8,0),gray,Arrow(4)); draw(o+(0,-3.2)--o+(0,3.2),gray,Arrow(4));
 label("$x$",o+(3.8,0),S,gray); label("$y$",o+(0,3.2),W,gray);}
pair o1=(0,0),o2=(6,0),o3=(0,-8),o4=(6,-8);
myax(o1); draw(shift(o1)*graph(f1,1/3,3.4),blue+linewidth(0.9));
myax(o2); draw(shift(o2)*graph(f2,1/3,3.4),blue+linewidth(0.9));
myax(o3); draw(shift(o3)*graph(g1,0,sqrt(3)),blue+linewidth(0.9)); draw(shift(o3)*graph(g2,0,3),red+dashed+linewidth(0.9));
myax(o4); draw(shift(o4)*graph(g3,0,sqrt(3)),blue+linewidth(0.9)); draw(shift(o4)*graph(g4,0,3),red+dashed+linewidth(0.9));
label("(a) $a>0,\ b<0$",o1+(1.8,-3.8)); label("(b) $a<0,\ b<0$",o2+(1.8,-3.8));
label("(c) $a>0,\ b>0$",o3+(1.8,-3.8)); label("(d) $a<0,\ b>0$",o4+(1.8,-3.8));
label("$y=x^{-1}$",o1+(1.2,0.9),NE,blue); label("$y=-x^{-1}$",o2+(1.2,-0.9),SE,blue);
label("$y=x^{2}$",o3+(sqrt(3),3),E,blue); label("$y=\sqrt{x}$",o3+(3,sqrt(3)),E,red);
label("$y=-x^{2}$",o4+(sqrt(3),-3),E,blue); label("$y=-\sqrt{x}$",o4+(3,-sqrt(3)),E,red);
label("\mbox{}",(min(currentpicture,user=true)+max(currentpicture,user=true))/2,black); // keeps the text after the figure black when LaTeX typesets the labels inline
\end{asy}
